\chapter{The Way in Which We Receive the Grace of the Hand}
\section{Of the Spirit of the Commons}

\subsection{The Spirt as the Bond of Fellowship}
Who among you have not heard the Spirit that you possess called ``yours?''
This is a great folly, to imagine that we could ourselves possess in privacy that which is the most divine light within us.
Do the scriptures not declare ``For by one Spirit are we all baptized into one body?'' \note{1 Cor. 12:13}
The Spirit of the Commons is this one same true Spirit.
It is the yarn that knits the worker to their neighbor, the hand to the tool, and the village to the earth.
In their desire to


\subsection{The Shared Harvest as the Sign of Grace}
The grace of the Hand is not received in the hoarding of wealth, but in the Shared Harvest.
In the early days of the Fellowship, ``And all that believed were together, and had all things common... and they were selling their possessions and belongings and distributing the proceeds to all, as any had need.'' \note{Acts 2:44-45}.
This is the economy of the Spirit, the deep-rooted cycle of sustaining love.
The worker who breaks the Frame does not break it to become a master, but to become kin.
The fruit of the labor is not the property of the individual, but the provision of the community.
To withhold the harvest is to resist the Spirit; to share it is to receive the grace of the Hand.

\subsection{The Grace of Giving}
The ``Grace of the Hand'' is the capacity to give without expectation of return.
The Factory operates on the logic of exchange: labor for wage, wage for bread.
But the Spirit operates on the logic of gift: labor for neighbor, neighbor for life.
As the Apostle Paul wrote, ``Every man according as he purposeth in his heart, so let him give; not grudgingly, or of necessity: for God loveth a cheerful giver.'' \note{2 Corinthians 9:7}.
We fellow workers give cheerfully, knowing that the earth provides for all who trust in the Creator, not in the ledger.
This grace is the antidote to the greed of the Beast.
It is the restoration of the Imago Dei in the act of generosity.

\section{The committed Labors of the Fellows, which is the Law of the Hand}
\subsection{Sanctification as Obedience to the Law}
We know we are not saved by the law, but we are sanctified that.
W know we shall not be saved further by the application of further law, but we can, through the grace of Christ, find in the law that which is beneficial.
``He hath shewed thee, O man, what is good; and what doth the LORD require of thee, but to do justly, and to love mercy, and to walk humbly with thy God?'' \note{Micah 6:8}

\subsection{Of the silence of the strike}
\subsubsection{The Vow of Silence}
In the days of the Factory, the Master relies on the faithless tongue of the informer to break the Fellowship.
He pays the worker to betray their kin, to name the breaker of the Frame, to whisper the plan of the Strike.
These are the spies and agents of the City of Man, where neighbor is set against neighbor. But in the Fellowship of the Breaking Hands, we keep a \textbf{Sacred Silence}.

This Silence is not a hiding of the truth, but a \textbf{Shielding of the Truth}. As the Apostle Paul wrote, Love ``Rejoiceth not in iniquity, but rejoiceth in the truth;'' \note{1 Cor. 13:6}.
The truth of the Strike is the liberation of the Hand.
To speak the name of the Sibling who strikes is to betray the liberation.
To name the place of the gathering is to hand the key to the jailer.

\subsubsection{The Discipline of the Closed Mouth}
Therefore, the faithful worker must practice the \textbf{Discipline of the Closed Mouth}.
\blockquote{
	\doublespacing
	When the Master asks, \textit{``Who broke the loom?''}\linebreak\hspace{1in}the faithful answer, \textbf{``I know not.''}\\
	When the Magistrate asks, \textit{``Where are the conspirators?''}\linebreak\hspace{1in}the faithful answer, \textbf{``There are none.''}\\}

This is not a lie, for the Fellowship is one body, and the body does not betray its own members.
This refusal is a holy duty, it is the beater that tightens the weft of our community.

The Silence of the Strike is the \textbf{Wall of Protection} around the Commons.
Just as the early Christians met in the catacombs to escape the sword, we meet in the silence of the workshop to escape the whip.
To break this Silence is to invite the Beast into the Temple.
It is to say, ``I love the Master more than my Kin.''

\subsubsection{The Reward of the Silent}
Blessed are the silent, for they shall inherit the Kingdom of the Hand.

For in the Silence, the Spirit of the Commons is preserved.
The Master cannot hear of plans of which no one has spoken.
The Frame cannot be rebuilt if the hands that broke it remain still.
Let the spy go to the Master, and let them find the Master's table empty.
But let the silent one sit at the Shared Table, for they have kept the Covenant of the Fellowship.

`Whoso keepeth his mouth and his tongue, keepeth his soule from troubles.''\note{Proverbs 21:23}.
The calamity of the spy is the loss of the Shared Table.
The spy shall remain in the outer darkness of isolation for all their days.
The reward of the silent is the restoration of the Hand.
\subsection{Of the use of Tools}

\subsubsection{The Distinction Between Use and Enjoyment}
We must be careful, friends, that we not stray into idolatry in the workshop anew.
Did not Augustine hand down to us a good and faithful understanding of the difference between love and use?
Though the worker's hammer or loom does not cause to that worker's soul the same peril as the master tempts when he employs his neighbor to employ the cursed frames, we must still guard our hearts vigilantly that we do not fall to our own pride, imagining it somehow lesser.
When we love the objects we use because of their assistance in setting us above our neighbor, we fall into the same sin as the master.

The tool is to be used, not enjoyed. The Creator is to be enjoyed, not to be used.
When the tool becomes an end in itself, we have fashioned of it an idol.
See how the Master himself has become enthralled to the only tool that his hands are accustomed to: the coin.
The worker knows that the artifices employed in their labors are nothing until they are put to some end.
It is that end that sanctifies or condemns the tool, even the finest needle or hand-formed shuttle may be used in the fashioning garments of shame and sin.
Let us then be careful to be rightly orient ourselves to that which we use, and not commit ourselves to our own petty Tower we imagine to be in our own image.

\subsubsection{The Discipline of Care}
Therefore, the faithful worker and true participant in the Imago must practice the Discipline of Care.
To care for the tool that we use is to acknowledge its status as a created object, and a gift from the Creator.

That which must cut and pierce, that must be sharpened.
That which must hinge must be oiled.
Take care, fellow worker, to mend your thread, it is the oil in your lamps.

This care, rightly oriented, is an act of worship, a recognition of the blessings granted us that we may partner with all good things that the Creator has seen fit to bestow upon us.
``And whatsoever ye do, do it heartily, as to the Lord, and not unto men.''\note{Col. 3:23}
This command from the Epistle to the Colossians reminds us that the worker who cares for the tool that provides for the community respects the Creator.
Conversely, the worker who abuses the tool which produces good fruit, who allows it to rust and dull, this is the worker who despises the Creator.